\documentclass[conference]{IEEEtran}
% \IEEEoverridecommandlockouts
% The preceding line is only needed to identify funding in the first footnote. If that is unneeded, please comment it out.
\usepackage{cite}
\usepackage{amsmath,amssymb,amsfonts}
\usepackage{algorithmic}
\usepackage{graphicx}
\usepackage{textcomp}
\usepackage{xcolor}
\usepackage{booktabs}
\usepackage{tabularx}
\usepackage{url}
\def\BibTeX{{\rm B\kern-.05em{\sc i\kern-.025em b}\kern-.08em
    T\kern-.1667em\lower.7ex\hbox{E}\kern-.125emX}}
\begin{document}
\raggedbottom

\title{Speech-Native Retrieval-Augmented Generation for Agricultural Advisory in Low-Resource Dialects \\
% {\footnotesize \textsuperscript{*}Note: Sub-titles are not captured in Xplore and
% should not be used}
% \thanks{Identify applicable funding agency here. If none, delete this.}
}
% Identify applicable funding agency here. If none, delete this.

\author{\IEEEauthorblockN{Darshan Pundlik Heble}
\IEEEauthorblockA{\textit{PG Scholar, Department of Computer Science} \\
\textit{Christ (Deemed to be University)}\\
Bangalore, India \\
darshanpundlik.heble@mca.christuniversity.in }
% \and
% \IEEEauthorblockN{2\textsuperscript{nd} Given Name Surname}
% \IEEEauthorblockA{\textit{dept. name of organization (of Aff.)} \\
% \textit{name of organization (of Aff.)}\\
% City, Country \\
% email address or ORCID}

}

\maketitle

\begin{abstract}
Smallholder farmers in low-literacy, dialect-rich regions face a persistent barrier to digital agricultural advisory: existing voice-enabled systems rely on a cascaded pipeline in which automatic speech recognition (ASR) first transcribes a spoken query before retrieval-augmented generation (RAG) produces an answer. Because domain-critical terminology — pesticide names, crop diseases, and pest names — is disproportionately subject to ASR substitution errors, and because farmers frequently use dialect-specific or colloquial terms absent from standard ASR training data, transcription errors propagate directly into retrieval failures and, ultimately, into incorrect agricultural advice. Recent work has shown that retrieval can be performed directly on speech embeddings without an intermediate transcription step, removing one source of cascade error; however, no existing system addresses the compounding effect of dialect variation and domain-critical rare vocabulary together, nor combines speech-native retrieval with an explicit dialect-to-scientific-entity mapping layer, confidence-gated escalation to a human expert, and fully offline operation on resource-constrained hardware. This paper proposes a framework that integrates these four components and outlines an experimental protocol — comparing ASR-cascaded, keyword-based, dense semantic, and hybrid retrieval under dialect variation — to evaluate whether this combination improves retrieval accuracy and answer reliability for agricultural advisory in low-resource dialect settings.
\end{abstract}

\renewcommand\IEEEkeywordsname{Keywords}
\begin{IEEEkeywords}
Agricultural advisory systems, dialect-aware retrieval, low-resource languages, speech-native RAG.
\end{IEEEkeywords}

\section{Introduction}
Agricultural extension — the transfer of evidence-based farming knowledge to smallholder farmers — remains a critical bottleneck in regions with low literacy and limited connectivity. RAG-based advisory systems such as Farmer.Chat \cite{farmerchat2024} and KrishokBondhu \cite{ameen2026krishokbondhu} show that voice-enabled, multilingual chatbots can reduce farmers' dependence on scarce extension agents, and systems such as SukhaRakshak AI \cite{iwmi2025sukharakshak} extend this to more than twenty Indian languages. All of them, however, first transcribe a spoken query via automatic speech recognition (ASR) before retrieval proceeds on the resulting text — a cascaded design vulnerable to a specific failure mode. ASR errors on Indian-language agricultural speech disproportionately affect domain-critical terms such as crop, chemical, and pest names \cite{benchmarkingasr2026}, and this degrades further under dialectal and code-mixed variation \cite{dialectmatters2026}. A mistranscribed pest or disease term propagates directly into retrieval, surfacing an irrelevant document and ultimately an incorrect recommendation.

One way to remove this failure mode is to retrieve directly from a speech representation, bypassing transcription entirely — a direction that has recently moved from proposal to validated, production-deployed technique. Amazon's SpeechRAG \cite{min2025speechrag} aligns a speech encoder with a frozen text-retrieval embedding space to retrieve audio passages without ASR, and Google's Speech-to-Retrieval (S2R) \cite{google2025s2r} performs an analogous mapping live in production for Voice Search. Neither system, however, has been evaluated under dialect variation or domain-critical rare-entity retrieval such as pesticide and pest names, and neither normalizes the many regional or colloquial terms a farmer may use for the same concept. Existing agricultural knowledge graphs \cite{wu2026cropgraphrag, kgeditorial2023} normalize only standardized, written terminology, leaving vernacular-to-scientific mapping for spoken queries an open problem.

A further limitation is that existing systems answer regardless of confidence, even though incorrect agricultural advice can damage a season's crop — an argument made explicitly for comparably high-stakes RAG deployments in clinical and financial domains \cite{dependablerag2026, abstinence2025}, though naive confidence signals are known to misfire in both directions \cite{intrygue2026} and require validation rather than blind trust. Most RAG research also implicitly assumes reliable connectivity and cloud-hosted inference \cite{compactllm2026}, which the rural settings these systems target frequently lack. Taken together, no existing system combines speech-native retrieval, dialect-to-entity mapping, confidence-gated escalation, and offline operation for agricultural advisory in low-resource dialects. This paper addresses that combination and poses the following research questions:

\begin{itemize}
    \item \textbf{RQ1:} Does speech-native retrieval retain its advantage over ASR-cascaded retrieval under dialect variation and domain-critical rare vocabulary?

    \item \textbf{RQ2:} Does a dialect-to-scientific-entity mapping layer improve retrieval accuracy for queries phrased in regional or colloquial terms?

    \item \textbf{RQ3:} Can the pipeline operate fully offline, on resource-constrained consumer hardware, at acceptable latency and accuracy?

    \item \textbf{RQ4:} Does hybrid (lexical + dense) retrieval combined with dialect mapping outperform either method alone on queries containing domain-critical rare entities?
\end{itemize}

The main contributions of this paper are:


\begin{itemize}
    \item Identifying a specific, underaddressed gap in speech-native agricultural retrieval---dialect variation combined with domain-critical rare vocabulary---not addressed by existing speech-native retrieval systems \cite{min2025speechrag, google2025s2r}.

    \item A dialect-to-scientific-entity mapping layer bootstrapped from low-cost, publicly available resources rather than new dialect speech collection.

    \item A confidence-gated RAG architecture that escalates low-confidence queries to a human expert.

    \item A system design and experimental protocol scoped to consumer-grade, GPU-constrained hardware, reserving cloud GPU resources only for one-time adapter training.
\end{itemize}


The remainder of this paper is organized as follows: Section II reviews related work; Section III presents the proposed architecture; Section IV describes the experimental setup; Section V presents results; and Section VI concludes.

\section{LITERATURE REVIEW}

% \subsection{Maintaining the Integrity of the Specifications}

\subsection{AI-Driven Agricultural Advisory Systems}

A growing body of work applies retrieval-augmented generation to agricultural advisory. Farmer.Chat \cite{farmerchat2024} is the most widely cited system, deployed across Kenya, India, Ethiopia, and Nigeria and engaging over 15,000 farmers across more than 300,000 queries; it supports more than 40 crops with multilingual, multimodal (text, voice note, image) input, ingesting structured and unstructured sources into a dynamic knowledge base. KrishokBondhu \cite{ameen2026krishokbondhu} is architecturally the closest system to the present proposal: a phone-call-based, fully voice-in/voice-out RAG pipeline for Bengali farmers in which OCR-digitized agricultural handbooks are vectorized, ASR transcribes the call, a vector database retrieves relevant passages, a small language model generates a grounded response, and text-to-speech delivers the answer; in pilot evaluation it produced high-quality answers for 72.7\% of queries and outperformed the KisanQRS benchmark \cite{rehman2023kisanqrs} by 44.7\% on a composite quality score. SukhaRakshak AI \cite{iwmi2025sukharakshak} integrates national language-AI infrastructure to deliver drought advisories in more than twenty Indian languages. AgroLLM \cite{agrollm2025} describes a simpler FAISS-based RAG chatbot for general agricultural question answering.

Three further systems are close enough in scope to merit individual discussion. A.A.H.A.R. \cite{dubey2025aahar} combines an offline, quantized LLM with an ASR component fine-tuned specifically for agricultural vocabulary and regional accents, reporting sub-2-second response times and high user satisfaction in its "Aahar" app; however, it remains a cascaded ASR-then-text architecture — its ASR is adapted to dialects, but retrieval still depends on the resulting transcript, and it includes neither speech-native retrieval nor confidence-gated escalation. Krishi Sathi \cite{vijayvargia2025krishisathi} is an intent-aware, multi-turn agricultural question-answering chatbot with ASR and TTS for accessibility in English and Hindi, reporting 97.53\% query response accuracy and sub-6-second response times, but follows the same ASR-then-retrieval design and does not address dialect variation specifically. A cross-lingual RAG case study for Bengali agricultural advisory \cite{hossain2026bengalirag} is presently text-only and already introduces a keyword-injection strategy to bridge colloquial farmer terminology with scientific nomenclature — a text-domain analog of this paper's proposed dialect-mapping layer — but its own discussion explicitly names ASR integration and dialect-aware normalization of spoken queries as future work it has not yet undertaken, independent confirmation from within the same problem space that this combination remains open.

None of the systems surveyed here treats ASR-induced retrieval failure on dialect-specific or domain-critical terminology as a first-class design problem to be solved by restructuring retrieval; the closest, A.A.H.A.R. and Krishi Sathi, address dialect variation only by adapting the ASR component itself.

\subsection{Speech Processing for Low-Resource Indian Languages and Dialects}

Benchmarking work on ASR for Indian-language agricultural speech \cite{benchmarkingasr2026} finds that standard word error rate (WER) under-penalizes errors on domain-critical terms, proposing an agriculture-weighted metric that shows the dominant ASR confusion pairs across languages involve crop, chemical, and pest or disease vocabulary specifically. A complementary study on cross-lingual ASR transfer for low-resource Indic dialects \cite{dialectmatters2026} quantifies how code-mixing and orthographic variation degrade transcription accuracy for dialect varieties such as Garhwali, building on national speech-data initiatives including Bhashini, Vakyansh, and IndicWav2Vec, and notes that no prior work had trained or evaluated an ASR model for that dialect at all. Dialect-specific speech corpora, such as a 250-hour Chhattisgarhi ASR dataset \cite{singh2023respin} produced under the RESPIN initiative and a dedicated Malayalam agricultural speech corpus, show that domain- and dialect-specific resources exist for some Indian languages but remain incomplete in coverage — most Indian dialects, including the one targeted in this paper, have no comparable resource yet.

\subsection{Speech-Native Retrieval}

SpeechRAG \cite{min2025speechrag} fine-tunes an adapter that projects speech-encoder representations into the embedding space of a frozen text retriever, enabling direct retrieval of audio passages from text or speech queries without an ASR step; retrieval experiments on spoken question-answering benchmarks show this approach matches or exceeds an ASR-then-text-retrieval cascade. Google's S2R \cite{google2025s2r} performs an analogous dual-encoder mapping in a production voice-search system and is accompanied by the Simple Voice Questions benchmark spanning seventeen languages. Underlying both approaches are self-supervised speech representation models such as wav2vec 2.0 \cite{baevski2020wav2vec} and HuBERT \cite{hsu2021hubert}, and the broader "textless NLP" literature \cite{spidr2025}, which has found that representations optimized for ASR accuracy are not necessarily optimal for downstream spoken-language tasks. Multilingual speech-translation foundation models such as SeamlessM4T \cite{barrault2023seamlessm4t} provide an additional source of pretrained multilingual speech encoders relevant to adapter-based approaches. Critically, neither SpeechRAG nor S2R has been evaluated under dialect variation or domain-critical rare-entity retrieval, and neither targets agriculture as an application domain — both are general-purpose systems whose published evaluations use standard, well-resourced language varieties.

\subsection{Retrieval Strategy: Keyword, Dense, and Hybrid Search}

Comparative studies of retrieval strategy \cite{sultania2024hybridsearch} consistently find that dense (semantic) retrieval outperforms lexical retrieval on paraphrased queries with little lexical overlap with the target document, while lexical methods such as BM25 retain an advantage on exact, rare, domain-specific entities — a category that includes pesticide and disease names directly relevant to agricultural queries. Hybrid retrieval, fusing lexical and dense signals, has been shown to outperform either method alone in domain-specific question answering. This is a practically important finding for the present proposal: a naive "speech embeddings beat keywords" framing would risk a result where keyword matching wins precisely on the safety-critical query category (pesticide and pest names), making hybrid retrieval — rather than a strict either/or comparison — the more defensible retrieval backbone.

\subsection{Dialect-Aware and Colloquial-to-Formal Knowledge Representation}

Agricultural knowledge graphs are a mature area of work, almost entirely built from standardized written text. Crop GraphRAG \cite{wu2026cropgraphrag} constructs a multi-relational knowledge graph spanning 28 crop categories, with retrieval performed jointly over graph subgraphs and text passages. The AgCNER dataset \cite{yao2024agcner} supports named-entity recognition for agricultural pests and diseases at scale. At the standards level, the EPPO ontology and related efforts \cite{kgeditorial2023} provide cross-system interoperable codes for plant and pest entities.

Closer to the present proposal's dialect-mapping layer, two recent works address the gap between informal farmer language and formal agricultural terminology — but neither in speech. Liu et al. \cite{liu2025colloquialpest} propose a soft-prompt-tuning method that extracts agricultural entities from colloquial, unstructured symptom descriptions using a pretrained agricultural language model, queries an agricultural knowledge graph for the matched entities, and classifies the pest or disease from the combined representation, reporting that conventional classifiers degrade substantially on farmers' natural, fuzzy language. The Bengali cross-lingual RAG case study \cite{hossain2026bengalirag} takes a simpler approach, injecting domain-specific keywords to bridge colloquial Bengali terms with scientific English nomenclature before retrieval. Both confirm that the colloquial-to-formal mapping problem is real, actively studied, and tractable — but both operate on written text, leaving the speech-domain version of this same mapping problem, applied to retrieval rather than classification, open.

\subsection{Confidence Estimation in RAG}

Dependable RAG work \cite{dependablerag2026} categorizes hallucination detection into closed-book methods, which rely on internal model signals, and open-book methods, which compare model behavior with and without retrieved context. Confidence-gated abstention has been argued for explicitly in high-stakes RAG deployments in clinical decision support and financial advisory, where the cost of a confident wrong answer exceeds the cost of abstaining \cite{abstinence2025}. However, naive entropy-based confidence signals have been shown to misfire in both directions, flagging correct, well-grounded answers as uncertain due to interactions between attention mechanisms and entropy computation \cite{intrygue2026}, indicating that a confidence-gating layer requires validation against labeled correct/incorrect examples rather than being assumed reliable by construction.

\subsection{Offline and Edge Deployment of Language Models}

Quantization and on-device deployment of large language models is an active systems research area. Reported figures for 4-bit quantized Llama-3 8B indicate usable but slow inference (2–3 tokens/second) on smartphones with at least 6 GB of RAM, with 8-bit variants requiring a discrete GPU with at least 8 GB of VRAM for comparable throughput \cite{compactllm2026}. Within agriculture specifically, A.A.H.A.R. \cite{dubey2025aahar} and the Bengali cross-lingual RAG case study \cite{hossain2026bengalirag} both report fully local, consumer-hardware deployments with sub-20-second end-to-end latency, demonstrating that offline agricultural RAG is already practically achievable for text-based systems; whether this holds once a speech-native retrieval and dialect-mapping layer are added is one of this paper's open questions (RQ3). A recent offline plant-disease detection system deployed in rural Tigray \cite{tigray2025plantdisease} demonstrates the broader deployment pattern — quantization, on-device runtime conversion, and dialect-localized interfaces — though for a vision rather than a speech modality.

\subsection{Synthesis}

Taken together, the systems reviewed above each address a subset of the requirements motivated in Section I, but none addresses all of them jointly. Farmer.Chat, KrishokBondhu, SukhaRakshak AI, A.A.H.A.R., and Krishi Sathi demonstrate that voice-enabled agricultural RAG is deployable at scale, but all rely on an ASR-then-retrieval cascade; A.A.H.A.R. and Krishi Sathi only adapt the ASR component to dialects rather than removing the dependence on transcription. SpeechRAG and S2R demonstrate that speech-native retrieval is viable and can outperform the cascade, but neither has been evaluated under dialect variation or domain-critical rare vocabulary, and neither targets agriculture. The colloquial-to-formal mapping work of Liu et al. and the Bengali RAG case study show that normalizing informal terminology is both necessary and actively studied, but only in text. Confidence-gated abstention and offline deployment are each independently established techniques, including within agricultural RAG specifically. The combination of all four — speech-native, dialect-aware, confidence-gated, offline agricultural retrieval — does not appear in any system located during this review, and one of the closest text-based systems in the same problem space explicitly names this combination as unaddressed future work of its own \cite{hossain2026bengalirag}.


\section{COMPARATIVE ANALYSIS OF PREVIOUS WORKS}



\section{Performance Analysis}


\subsection{Impact of Knowledge Graphs on Retrieval Quality}


\subsection{Efficacy of Multi-Agent Filtering}


\subsection{Domain-Specific Accuracy: The Case of Scientific Data}


\section{Conclusion}

Conclusion will be here


% \subsection{Abbreviations and Acronyms}\label{AA}
% Define abbreviations and acronyms the first time they are used in the text, 


% \subsection{Units}
% \begin{itemize}
% \item Use either SI
% \item Avoid 
% \end{itemize}

% \subsection{Equations}
% Number 
% \begin{equation}
% a+b=\gamma\label{eq}
% \end{equation}


% \subsection{Figures and Tables}
% \paragraph{Po 
% ``Fig.~\ref{fig}'', even at the beginning of a sentence.

% \begin{table}[htbp]
% \caption{Table Type Styles}
% \begin{center}
% \begin{tabular}{|c|c|c|c|}
% \hline
% \textbf{Table}&\multicolumn{3}{|c|}{\textbf{Table Column Head}} \\
% \cline{2-4} 
% \textbf{Head} & \textbf{\textit{Table column subhead}}& \textbf{\textit{Subhead}}& \textbf{\textit{Subhead}} \\
% \hline
% copy& More table copy$^{\mathrm{a}}$& &  \\
% \hline
% \multicolumn{4}{l}{$^{\mathrm{a}}$Sample of a Table footnote.}
% \end{tabular}
% \label{tab1}
% \end{center}
% \end{table}

% \begin{figure}[htbp]
% \centerline{\includegraphics{fig1.png}}
% \caption{Example of a figure caption.}
% \label{fig}
% \end{figure}

% Figure Labels: Use 8 point Times New Roman for Figure labels. Use words 
% rather than symbols or abbreviations when writing Figure axis labels to 
% avoid confusing the reader. As an example, write the quantity 
% ``Magnetization'', or ``Magnetization, M'', not just ``M''. If including 
% units in the label, present them within parentheses. Do not label axes only 
% with units. In the example, write ``Magnetization (A/m)'' or ``Magnetization 
% \{A[m(1)]\}'', not just ``A/m''. Do not label axes with a ratio of 
% quantities and units. For example, write ``Temperature (K)'', not 
% ``Temperature/K''.


% \section*{Acknowledgment}

% The preferred spelling of the word ``acknowledgment'' in America is without 
% an ``e'' after the ``g''. Avoid the stilted expression ``one of us (R. B. 
% G.) thanks $\ldots$''. Instead, try ``R. B. G. thanks$\ldots$''. Put sponsor 
% acknowledgments in the unnumbered footnote on the first page.

\bibliographystyle{IEEEtran}
\bibliography{reference}

\end{document}
